\section{Conclusion}

We formulate cross-environment interpretability of MARL-GPT as circuit tracing rather than single-layer representation analysis. Two independent full-path CLTs expose sparse computations for the actor and critic, and local attribution graphs connect structured input tokens to legal-action contrasts and action-conditioned value predictions through residual and attention-OV paths. The central empirical burden is not producing an attractive graph: it is showing that the replacement remains faithful, reconstruction errors do not carry the explanation, graph motifs recur under matched conditions, and interventions have the predicted effects in the original policy and in fresh rollouts. Meeting those gates would support a mechanistic account of cross-environment generalization and a principled route to lightweight policy steering. Until then, the framework defines a falsifiable experiment rather than a claim that shared mechanisms have been found.
